\newcommand{\myrule}{\vskip0.15in\hrule\vskip0.10in}

\begin{figure*}

{\small

\begin{minipage}[t]{0.31\textwidth}

\myrule

{\begin{verbatim}
%0:i16 = var
%1:i16 = or %0, 5:i16
%2:i16 = and %1, 4:i16
infer %2
\end{verbatim}}
$\Rightarrow$
{\begin{verbatim}
result 4:i16
\end{verbatim}}

This example, and others like it, motivated us to implement a new
\emph{bit estimator} analysis in the Visual C++ compiler that identifies bits that are
provably zero or one, making this optimization and many related ones
easy to implement.

\myrule

{\begin{verbatim}
%0:i32 = var
%1:i32 = xor %0, 4294967295:i32
%2:i32 = and %1, 8:i32
infer %2
\end{verbatim}}
$\Rightarrow$
{\begin{verbatim}
%3:i32 = and 8:i32, %0
%4:i32 = xor 8:i32, %3
result %4
\end{verbatim}}

Instead of inverting the entire word and then isolating bit~3, we can
first isolate the bit and then invert it.
%
The resulting code contains a 1-byte immediate value instead of a
4-byte one.
%
(\tool{}'s cost function is not currently clever enough to capture
this fact---it was a matter of luck that this optimization was
synthesized.)

\myrule{
\begin{verbatim}
%0:i8 = var
%1:i8 = lshr %0, 3:i8
%2:i1 = eq %1, 0:i8
infer %2
\end{verbatim}}
$\Rightarrow$
{\begin{verbatim}
%3:i1 = ult %0, 8:i8
result %3
\end{verbatim}}

A value shifted right 3 times can be zero only if it is smaller than $2^3 = 8$.

\end{minipage}
%
\hspace{0.035\textwidth}
%
\begin{minipage}[t]{0.31\textwidth}
\myrule
{\begin{verbatim}
%0:i64 = var
%1:i64 = and %0, 31:i64
%2:i64 = add %1, 1:i64
%3:i1 = ult 32:i64, %2
infer %3
\end{verbatim}}
$\Rightarrow$
{\begin{verbatim}
result 0:i1
\end{verbatim}}

Adding 1 to a value in range $0..31$ cannot produce a result larger
than 32. The bit estimator can be used to implement the family of
optimizations suggested by this example.

\myrule
{\begin{verbatim}
%0:i64 = var
%1:i64 = mul %0, 2654435761:i64
%2:i64 = and %1, 1:i64
%3:i1 = ne %2, 0:i64
infer %3
\end{verbatim}}
$\Rightarrow$
{\begin{verbatim}
%4:i1 = trunc %0
result %4
\end{verbatim}}

Since multiplying two odd numbers always yields an odd number, the
result can be odd only if \%0 is odd to begin with.

\myrule
{\begin{verbatim}
%0:i32 = var
%1:i64 = zext %0
%2:i64 = udiv -1:i64, %1
%3:i1 = ule 4:i64, %2
infer %3
\end{verbatim}}
$\Rightarrow$
{\begin{verbatim}
result 1:i1
\end{verbatim}}

Dividing the largest unsigned 64-bit number by an unsigned 32-bit
number cannot produce a value that is smaller than $2^{32} - 1$.
%
We used the bit estimator to prove a lower bound on the value of \%2,
making is possible to fold the comparison to true.

\end{minipage}
%
\hspace{0.035\textwidth}
%
\begin{minipage}[t]{0.31\textwidth}

\myrule
{\begin{verbatim}
%0:i32 = var
%1:i1 = eq %0, 0:i32
%2:i32 = select %1, 2:i32, 1:i32
infer %2
\end{verbatim}}
$\Rightarrow$
{\begin{verbatim}
%3:i1 = ult %0, 1:i32
%4:i32 = zext %3
%5:i32 = shl 1:i32, %4
result %5
\end{verbatim}}

\tool{} excels at replacing \texttt{select} instructions, which
usually become CMOV instructions on x86/x86-64, with logical
operations.

\myrule
{\begin{verbatim}
%0:i1 = var
%1:i8 = select %0, 1:i8, 0:i8
%2:i1 = ne %1, 1:i8
%3:i32 = select %2, 40:i32, 20:i32
infer %3
\end{verbatim}}
$\Rightarrow$
{\begin{verbatim}
%4:i32 = zext %0
%5:i32 = ashr 40:i32, %4
result %5
\end{verbatim}}

Another \texttt{select} removal.

\myrule
{\begin{verbatim}
%0:i32 = var
%1:i32 = and %0, 131071:i32
%2:i32 = mul %1, 65536:i32
infer %2
\end{verbatim}}
$\Rightarrow$
{\begin{verbatim}
%3:i32 = shl %0, 16:i32
result %3
\end{verbatim}}

\tool{} is good at removing useless instructions.

\myrule
{\begin{verbatim}
%0:i64 = var
%1:i64 = add %0, -10445360463872:i64
%2:i64 = and %1, 4095:i64
infer %2
\end{verbatim}}
$\Rightarrow$
{\begin{verbatim}
%3:i64 = and 4095:i64, %0
result %3
\end{verbatim}}

Another useless instruction removed.

\end{minipage}

}

\vskip 0.3in

\caption{Representative optimizations suggested by \tool{} for LHSs
  extracted from the Windows kernel. We implemented generalized
  versions of these, and others, in the Visual C++ compiler, adding a
  total of 40 new optimizations.}
\label{fig:windows}
\end{figure*}



\iffalse
Useless instructions
Instructions that don’t affect the results of the entire expression

{\begin{verbatim}
%0:i64 = var
%1:i64 = and %0, -8:i64
%2:i64 = lshr %1, 12:i64
%3:i64 = add %1, 7:i64
%4:i64 = lshr %3, 12:i64
%5:i1 = eq %2, %4
infer %5
\end{verbatim}}
$\Rightarrow$
{\begin{verbatim}
result 1:i1
\end{verbatim}}

The addition with 7 (marked in bold) is useless, since:

None of the bits from \%2 after bit 3 are affected, since the and
of \%1 with -8 resets the low 3 bits.

The resulting value is shifted right by 12 bits, discarding the 3 bits just added.

After the addition is removed, value numbering discovers that \%2
and \%4 are equivalent, turning the entire expression into the result
1. This pattern originated from the following code sequence found in
the Windows Kernel, which is quite common when aligning/rounding up
pointers and memory allocation size values.

\begin{verbatim}
StackBuffer = ALIGN_UP_POINTER_BY(StackBuffer, HV_CALL_ALIGNMENT);
// If the stack buffer crosses a page, align it.
if (((ULONG_PTR)StackBuffer >> PAGE_SHIFT) !=
    (((ULONG_PTR)StackBuffer + ElementSize - 1) >> PAGE_SHIFT)) {
  StackBuffer = PAGE_ALIGN((ULONG_PTR)StackBuffer + ElementSize - 1);
}
\end{verbatim}

\fi

